\documentclass[12pt]{article}

\usepackage{sbc-template}
\usepackage{graphicx,url}
\usepackage[utf8]{inputenc}
\usepackage[brazil]{babel}
\usepackage[latin1]{inputenc}  

     
\sloppy

\title{Avaliação de Usabilidade do Sistema Gestão de Frequência: \\ Uma Abordagem Baseada nas Heurísticas de Nielsen}

\author{Hygor Albert Fernandes Marques\inst{1}}

\address{Instituto Federal de Santa Catarina (IFC) -- Campus Videira \\
  Videira -- SC -- Brasil
  \email{hygor.marques@estudantes.ifc.edu.br}
}

\begin{document} 

\maketitle

\begin{abstract}
  fazer depoios nao esquecer
\end{abstract}
     
\begin{resumo} 
  fazer depoios nao esquecer
\end{resumo}


\section{Introdução}

O monitoramento da frequência estudantil é um dos pilares para a mitigação da evasão escolar e para o suporte pedagógico preventivo. 
Entretanto, sistemas de gestão acadêmica consolidados, como o SIGAA (Sistema Integrado de Gestão de Atividades Acadêmicas), 
frequentemente priorizam a funcionalidade de registro em detrimento da visualização analítica dos dados. Para o gestor escolar, 
o acesso a tabelas brutas sem o devido processamento de informações gera uma alta carga cognitiva, 
dificultando a identificação de padrões de absenteísmo que ex//igem intervenção imediata.

A usabilidade de um Ambiente Virtual de Aprendizagem (AVA) ou de um sistema de gestão não deve ser vista apenas como um atributo estético, 
mas como uma ferramenta de eficiência operacional. Conforme apontado por Antonino e Freire (2017), 
falhas na interface e na forma como o sistema auxilia o usuário na correção e interpretação de dados podem comprometer a eficácia do processo de gestão. 
Diante desse cenário, surge a necessidade de interfaces que não apenas armazenem dados, mas que os transformem em conhecimento estratégico.

Este artigo descreve o desenvolvimento e a avaliação de usabilidade do sistema "Gestão de Frequência", 
uma plataforma web projetada para o monitoramento inteligente de faltas. O objetivo central deste trabalho é 
demonstrar como a aplicação das 10 Heurísticas de Nielsen pode evoluir uma interface puramente funcional para um dashboard de 
gestão diagnóstica. Através de um design minimalista e focado no usuário, o sistema busca responder a questões críticas para a gestão escolar: 
quais alunos estão em risco, em quais disciplinas as faltas se concentram e quais dias da semana apresentam maior índice de evasão.

\end{document}
